\documentclass[11pt]{article}

\usepackage[margin=1in]{geometry}
\usepackage[T1]{fontenc}
\usepackage{lmodern}
\usepackage{amsmath, amssymb}
\usepackage{setspace}
\usepackage[most]{tcolorbox}
\usepackage{xcolor}
\usepackage{fancyhdr}

\tcbset{
  enhanced,
  boxrule=0pt,
  frame hidden,
  breakable
}

\pagestyle{fancy}
\fancyhf{}
\setlength{\headheight}{52pt}
\renewcommand{\headrulewidth}{0pt}

\newcommand{\Course}{Math 4610 Spring 2026}
\newcommand{\Instructor}{Discussion 3}
\newcommand{\AssignmentType}{Discussion}
\newcommand{\AssignmentNumber}{3}
\newcommand{\Contributors}{
Marks, Stephen\\
}

\fancyhead[L]{\Course\\ \Instructor}
\fancyhead[R]{\Contributors}

\newcommand{\ProblemDivider}[1]{%
  \vspace{1.0em}
  \noindent\textbf{#1}\par
  \vspace{0.35em}
}

\definecolor{problemBack}{RGB}{235,242,250}
\definecolor{problemFrame}{RGB}{70,110,160}
\definecolor{exerciseGray}{gray}{0.96}
\definecolor{exerciseGrayFrame}{gray}{0.45}

\newtcolorbox{problemBox}{
  colback=exerciseGray,
  colframe=exerciseGrayFrame,
  arc=2mm,
  left=8pt,right=8pt,top=8pt,bottom=8pt
}

\newtcolorbox{solutionBox}{
  colback=white,
  colframe=white,
  arc=0mm,
  left=0pt,right=0pt,top=0pt,bottom=0pt
}

\newcommand{\ProblemAnnounce}{\noindent}
\newcommand{\SolutionAnnounce}{\noindent\textbf{Solution.}\par\medskip}
\newcommand{\Exercise}[1]{\textbf{Exercise~#1.}\;}

\newcommand{\R}{\mathbb{R}}
\newcommand{\C}{\mathbb{C}}
\newcommand{\F}{\mathbb{F}}

\begin{document}

\begin{center}
  {\Large \textbf{\AssignmentType~\AssignmentNumber}}
\end{center}
\vspace{0.75em}

\ProblemDivider{1}

\begin{problemBox}
\ProblemAnnounce
\Exercise{5.A.17}
Show that the forward shift operator \(T \in \mathcal{L}(\F^{\infty})\) defined by
\[
T(z_1, z_2, \ldots) = (0, z_1, z_2, \ldots)
\]
has no eigenvalues.
\end{problemBox}

\begin{solutionBox}
\SolutionAnnounce
\doublespacing
Suppose \(\lambda\) is an eigenvalue of \(T\) with eigenvector \(\mathbf{z} = (z_1, z_2, \dots) \neq \mathbf{0}\). Then
\[
T\mathbf{z} = \lambda \mathbf{z}.
\]
So, 
\[
(0, z_1, z_2, \dots) = (\lambda z_1, \lambda z_2, \lambda z_3, \dots).
\]
Then immediately, \(0 = \lambda z_1\).

\begin{itemize}
    \item If \(\lambda = 0\), then the remaining equations give \(z_1 = 0\), \(z_2 = 0\), \(z_3 = 0\), and so on. So \(\mathbf{z} = \mathbf{0}\), which is a contradiction. So \(\lambda = 0\) is not an eigenvalue.
    \item If \(\lambda \neq 0\), then from \(0 = \lambda z_1\) we get \(z_1 = 0\). The second equation gives \(z_1 = \lambda z_2\), so \(0 = \lambda z_2\) and hence \(z_2 = 0\). Then all \(z_k = 0\), again contradicting \(\mathbf{z} \neq \mathbf{0}\).
\end{itemize}
Therefore, no \(\lambda \in \F\) can be an eigenvalue of \(T\). Thus \(T\) has no eigenvalues.
\end{solutionBox}

\newpage
\ProblemDivider{2}

\begin{problemBox}
\ProblemAnnounce
\Exercise{5.A.18}
Define the backward shift operator \(S \in \mathcal{L}(\F^{\infty})\) by
\[
S(z_1, z_2, \dots) = (z_2, z_3, \dots).
\]
\begin{enumerate}
    \item Show that every element of \(\F\) is an eigenvalue of \(S\).
    \item Find all the eigenvectors of \(S\).
\end{enumerate}
\end{problemBox}

\begin{solutionBox}
\SolutionAnnounce
\doublespacing

\begin{enumerate}
    \item Let \(\lambda \in \F\). Consider a nonzero \(\mathbf{z} = (z_1, z_2, \dots)\) such that
    \[
    S\mathbf{z} = \lambda \mathbf{z}.
    \]
    This means
    \[
    (z_2, z_3, \dots) = (\lambda z_1, \lambda z_2, \lambda z_3, \dots).
    \]
    Comparing components gives: 
    \[
    z_{k+1} = \lambda z_k \quad \text{for all } k \ge 1.
    \]
    Choose \(z_1 = 1\). Then: 
    \[
    z_2 = \lambda, \quad z_3 = \lambda^2, \quad z_4 = \lambda^3, \dots
    \]
    Thus the vector
    \[
    \mathbf{z}_\lambda = (1, \lambda, \lambda^2, \lambda^3, \dots)
    \]
    is nonzero and satisfies \(S\mathbf{z}_\lambda = \lambda \mathbf{z}_\lambda\). So every \(\lambda \in \F\) is an eigenvalue of \(S\).

    \item The eigenspace for \(\lambda\) is
    \[
    E_\lambda = \{ c \cdot (1, \lambda, \lambda^2, \dots) \mid c \in \F,\ c \neq 0 \}.
    \]
    Therefore, the set of all eigenvectors of \(S\) is the union over all \(\lambda \in \F\) of \(E_\lambda\).
\end{enumerate}

\end{solutionBox}

\end{document}
